\documentclass{article}

\usepackage{bm}
\usepackage{ctex}
\usepackage{tikz}
\usepackage{float}
\usepackage{xcolor}
\usepackage{setspace}
\usepackage{datetime}
\usepackage{geometry}
\usepackage{graphicx}
\usepackage{enumitem}
\usepackage{tcolorbox}
\usepackage{nicematrix}
\usepackage{amsmath, amssymb, amsfonts, mathrsfs}
\usepackage[fontsize = 12pt]{fontsize}

\renewcommand{\thesubsection}{(\alph{subsection})}
\newcommand{\diff}[2]{\dfrac{\mathrm{d}#1}{\mathrm{d}#2}}
\newcommand{\diffn}[3]{\dfrac{\mathrm{d}^{#3}#1}{\mathrm{d}#2^{#3}}}
\newcommand{\piff}[2]{\dfrac{\partial{#1}}{\partial{#2}}}
\newcommand{\eval}[2]{\left.#1\right|_{#2}}
\newcommand{\e}{\mathrm{e}}
\renewcommand{\d}{\mathrm{d}}
\newcommand{\barline}[1]{\overline{#1\rule{0pt}{1.5ex}}}

\geometry{
    a4paper,
    left = 3.18cm,
    right = 3.18cm,
    top = 2.54cm,
    bottom = 2.54cm
}

\setitemize[1]{
    itemsep = 0pt,
    partopsep = 0pt,
    parsep = \parskip,
    topsep = 5pt
}

\title{应用统计：第一次作业}
\author{叶炳辉 \quad 202528017729004 \quad 人工智能学院}
\date{\currenttime,\ \today}



\begin{document}

\maketitle



\section{Q1}
\textit{\textcolor{gray}{从某鱼池中取 100 条鱼，做上记号后放在该鱼池中。先从该池中任意捉来 40 条鱼，发现其中 2 条有记号，问该鱼池中大约有多少条鱼？}}

设该鱼池中大约有总鱼数为 $N$ 条，根据题意：

\begin{itemize}
    \item 整个鱼池中有记号的鱼的数量为 100 条。
    \item 样本的总数量为 40 条。
    \item 样本中有记号的数量为 2 条。
\end{itemize}

由此可以列出等式：
\[ \dfrac{100}{N}=\dfrac{2}{40}\implies N=2000 \]

因此，该鱼池中大约有 2000 条鱼。



\section{Q3}
\textit{\textcolor{gray}{设某种动物由出生算起，活到 20 年以上的概率为 0.8，活到 25 年以上的概率为 0.4。如果现在有一只活了 20 年的这种动物，问它能活到 25 年以上的概率是多少？}}

设事件 $A$ 为“该动物活到 20 年以上”，事件 $B$ 为“该动物活到 25 年以上”，则：
\[ \Pr(A)=0.8\qquad \Pr(B)=0.4 \]

因为一只动物活到 25 年以上，就意味着它必然已经活到了 20 年以上，所以事件 $B$ 是事件 $A$ 的子集，即 $B\subseteq A$，则事件 $A$ 和事件 $B$ 同时发生的概率就等于事件 $B$ 发生的概率：
\[ \Pr(AB)=\Pr(B)=0.4 \]

在已知该动物已经活了 20 年（即事件 $A$ 已经发生）的条件下，它能活到 25 年以上（即事件 $B$ 发生）的概率，根据条件概率公式：
\[ \Pr(B\mid A)=\dfrac{\Pr(BA)}{\Pr(A)}=\dfrac{0.4}{0.8}=50\% \]

综上，它能活到 25 年以上的概率是 50\%。



\section{Q5}
\textit{\textcolor{gray}{已知一批玉米种子的出苗率为 0.9，现每穴中两粒。}}

已知每粒种子的出苗率是相互独立的，且出苗的概率为 $p=0.9$，不出苗的概率为 $q=1-p=0.1$，每穴播种 $n=2$ 粒种子，如果设每穴中出苗的种子数为随机变量 $X$，则 $X$ 服从参数为 $(2,\ 0.9)$ 的二项分布，即 $X\sim B(2,\ 0.9)$。

根据二项分布的概率公式：
\[ \Pr(X=k)=\binom{n}{k}p^{k}q^{n-k} \]



\subsection{两粒都出苗的概率}
\[ \Pr(X=2)=\binom{2}{2}\cdot0.9^{2}\times0.1^{2-2}=0.81 \]



\subsection{一粒出苗另一粒不出苗的概率}
\[ \Pr(X=1)=\binom{2}{1}\cdot0.9^{1}\times0.1^{2-1}=0.18 \]



\subsection{两粒都不出苗的概率}
\[ \Pr(X=0)=\binom{2}{0}\cdot0.9^{0}\times0.1^{2-0}=0.01 \]



\end{document}
